
\section{Generative Models}
\label{sec:generative}

% ---------------------------------------------------------------
% Sub-topics to cover:
%
% 1. Variational Autoencoders (VAEs)
%    - Encoder q_\phi(z|x), decoder p_\theta(x|z).
%    - ELBO objective; reparameterization trick.
%    - Limitation: blurry samples due to pixel-wise MSE.
%    - HEP use: fast calorimeter simulation (e.g., CaloVAE).
%
% 2. Generative Adversarial Networks (GANs)
%    - Generator vs. discriminator; minimax game.
%    - Training instability; mode collapse.
%    - HEP use: CaloGAN, jet generation.
%
% 3. Normalizing flows
%    - Bijective f: x <-> z with tractable Jacobian.
%    - Change of variables formula: log p(x) = log p(z) + log|det J|.
%    - Coupling layers (RealNVP, Glow); autoregressive flows (MAF).
%    - HEP use: fast simulation (SurVAE, NF-based shower generators),
%      inference (nflows for posterior estimation).
%    - Key property: exact density evaluation — useful for reweighting.
%
% 4. Motivation for moving to diffusion
%    - Flows require bijective transformations → architectural
%      constraints.  GANs are unstable.  VAEs are blurry.
%    - Diffusion models sidestep these issues at the cost of
%      sequential sampling.  Forward pointer to \cref{sec:diffusion}.
